% FIGURE NEEDED: solution sketch of the eclipse light curve labelling the
% eclipse duration t_e (A to D) and total-eclipse duration t_t (B to C)
% (2008_th_sol.pdf, page 9).

Binary period: $P = 30$ days $= 0.0821$ years.

Time of eclipse: $t_e = 8$ hours $= 0.0009$ years.

Time of total eclipse: $t_t = 1$ hour 18 minutes $= 0.0001$ years.

Radial velocity of primary star: $V_{r1} = 30\,\mathrm{km/s}
= 6.3285$\,AU/year.

The velocity of secondary star: $V_{r2} = 40\,\mathrm{km/s}
= 8.4380$\,AU/year.

For circular orbit
\[ V_r = \frac{2\pi a}{P} \quad\Longrightarrow\quad
   a = \frac{P V_r}{2\pi} \quad (a = \text{radius of orbit}) \]

For primary star: \hfill(10\%)
\[ a_1 = \frac{P V_{r1}}{2\pi}
   = \frac{(0.0821)(6.3285)}{2\pi}
   = 0.0827\ \mathrm{AU} = 1.24 \times 10^7\ \mathrm{km} \]

For secondary star: \hfill(10\%)
\[ a_2 = \frac{P V_{r2}}{2\pi}
   = \frac{(0.0821)(8.4380)}{2\pi}
   = 0.1103\ \mathrm{AU} = 1.65 \times 10^7\ \mathrm{km} \]
\[ a = a_1 + a_2 = 0.0827 + 0.1103
   = 0.1930\ \mathrm{AU} = 2.89 \times 10^7\ \mathrm{km} \]

The radius of the primary star can be determined from the equation:
\hfill(10\%)
\[ \frac{(t_e + t_t)}{P} = \frac{4R_1}{2\pi a}
   \qquad\text{or}\qquad
   R_1 = \frac{2\pi a}{4P}(t_e + t_t)
   = \frac{2\pi(0.1930)}{4(0.0821)}(0.0009 + 0.0001) \]
\[ = 0.0039\ \mathrm{AU} = 5.86 \times 10^5\ \mathrm{km}
   = 0.84\ R_\odot \]
\hfill(20\%)

The radius of the secondary star can be determined from the equation:
\hfill(10\%)
\begin{align*}
\frac{(t_e - t_t)}{P} = \frac{4R_2}{2\pi a}
  \qquad\text{or}\qquad
  R_2 &= \frac{2\pi a}{4P}(t_e - t_t)
  && \text{(10\%)} \\
  &= \frac{2\pi(0.1930)}{4(0.0821)}(0.0009 - 0.0001) \\
  &= 0.0028\ \mathrm{AU} = 4.22 \times 10^5\ \mathrm{km} = 0.61\ R_\odot
  && \text{(15\%)}
\end{align*}

From the Kepler's third law:
\[ \frac{a^3 \sin^3 i}{P^2} = \left(\mathcal{M}_1 + \mathcal{M}_2\right)
   \qquad\text{and}\qquad
   \frac{\mathcal{M}_1}{\mathcal{M}_2} = \frac{a_2}{a_1} \]

Since $i = 90^\circ$, $\sin i = \sin 90^\circ = 1$, and then
\[ \mathcal{M}_1 = \frac{a^3}{P^2(1 + a_1/a_2)}
   = \frac{(0.1930)^3}{(0.0821)^2(1 + 0.0827/0.1103)}
   = 0.61\ \mathcal{M}_\odot \]
\hfill(15\%)
\[ \mathcal{M}_2 = \frac{a^3}{P^2(1 + a_2/a_1)}
   = \frac{(0.1930)^3}{(0.0821)^2(1 + 0.1103/0.0827)}
   = 0.46\ \mathcal{M}_\odot \]
\hfill(10\%)
